% Proposed SI prose. No frozen benchmark or manuscript file was modified by this audit.
\subsection*{Current identifier consistency and historical metadata dependencies}
We checked all 1,035 frozen benchmark SMILES against the current PubChem properties of their 1,030 distinct CIDs on 3 October 2026. All 26 batched property requests returned the requested records. Standard InChIKeys were calculated with RDKit 2025.09.4. We compared full keys, their connectivity blocks, and, where these differed, keys after RDKit FragmentParent processing. Responses, request times, hashes and row-level comparisons are archived in \path{supplementary_audits/identity_sources/}. These current responses do not reconstruct the 2024 source snapshot or establish the identity of historical test materials.

Full InChIKeys agreed for 875 rows. Another 127 rows shared the connectivity block but differed in the full key. Of the remaining 33 rows, six were related under fragment-parent processing and 27 lacked that correspondence. These categories identify records for review; they are not estimates of erroneous toxicity labels. In particular, fragment correspondence cannot establish equivalence between salts, mixtures or test materials. Before examining the resulting model sensitivity, we froze the 33-row \texttt{current\_identifier\_flagged} list using identifier comparisons alone. The sensitivity excludes those identities from both training and testing within each original split, with identical retained members at all three endpoints. It preserves the frozen SMILES, labels and official split assignments for retained identities. The 127 connectivity-matched rows were not automatically excluded from this non-chiral ECFP sensitivity.

Targeted review of five repeated-CID groups distinguished different cases. The current CIDs recorded for silthiofam, diflumetorim and phthalide identify different structures; CAS queries returned CIDs 9881821, 15390131 and 33791, respectively, whose full InChIKeys agree with the corresponding frozen SMILES. Chlormequat requires a cation/salt distinction: CAS 7003-89-6 returns CID 13837 for the cation, whereas CID 13836 matches the chloride-salt row. Metiram is distinct from zineb despite their shared recorded CID; the current PubChem reference-substance record for metiram provides no discrete CID, so no replacement CID was assigned. Metominostrobin shares connectivity with its current CID but lacks the latter's specified oxime geometry in the frozen SMILES. The complete ten-row ledger records these distinctions without replacing historical structures or labels.

The pinned upstream code obtains the year field from CID-linked PubChem literature dates, returning the first acceptable dated entry among five retrieved records after excluding three named problematic titles. When that routine returns no usable date, it takes the CID record's creation date. It then extracts the year and sorts the combined dataset for the Time split. Thus year is a literature-or-record-creation proxy, rather than a toxicity-test date or a uniformly verified first-publication year. The historical per-row date responses and branch choices are unavailable. Attempts to replay the legacy literature endpoint for the nine stored or candidate CIDs returned HTTP 404 HTML responses; these failures were not interpreted as absence of publications and did not trigger substitute year assignments. No specific correction to the frozen years or Time membership was established.

The metadata dependency differs by source. ECOTOX use flags were parsed from CID-linked PubChem agrochemical information; PPDB/BPDB flags were parsed from their provider HTML. The three named current CID mismatches above occur in PPDB rows, so their CID discrepancies do not by themselves establish incorrect use flags. The two ECOTOX rows within these repeated-CID groups have current full-key agreement. Downstream year and use effects remain unresolved unless separately supported; absence of an established error is not proof of correctness.

\subsection*{Additional source checks and their effects on the frozen audit}
Replaying the upstream exclusion-log logic distinguished logged rows from rejected route groups. The log contains 745 observations across 72 CAS identities because it selects all observations sharing a CAS with any unresolved route. The actual 11-cut-off filter rejects 79 route groups containing 536 observations; another 209 observations in 75 retained route groups appear in the same log. This logging behaviour does not alter the route-level filter itself. The row-level audit preserves both flags.

In the official Rothamsted report, Table 2 (printed page 56; PDF page 3) reports benomyl oral LD$_{50}$ as $>10$~\ugbee, whereas the historical ECOTOX row reports $<10$~\ugbee. The archived PDF has SHA256 \texttt{270ae180f6a58f7c03133ff6b226d5b9ca2c88dd7174ece6b4d8c80e6c039ae1}. We treat this as a document-level discrepancy: complete one-to-one trial matching and independent adjudication by two named experts have not been established. A conditional propagation check, accepting the original-table mapping while retaining the upstream-eligible route set, changes the primary-rule labels at 100/11/1 from positive/positive/unresolved to unresolved/unresolved/negative. The keep policy nevertheless retains the same final labels, and the exclude policy removes this identity under both versions. This local check supports unchanged membership and final labels for those two primary policies, not a general assertion that all audit counts or sensitivity scenarios are unaffected. The proposed source interpretation was not written into the frozen records.
